\documentclass[letterpaper]{article} % DO NOT CHANGE THIS
\usepackage[submission]{aaai2027}  % DO NOT CHANGE THIS

% AAAI-compatible packages.
\usepackage[hyphens]{url}
\usepackage{graphicx}
\usepackage{natbib}
\usepackage{caption}
\usepackage{booktabs}
\usepackage{longtable}
\usepackage{array}
\usepackage{amsmath}
\usepackage{amssymb}
\usepackage{algorithm}
\usepackage{algorithmic}
\usepackage{multicol}
\usepackage{xr}
\setcounter{topnumber}{5}
\setcounter{dbltopnumber}{5}
\setcounter{totalnumber}{10}
\renewcommand{\topfraction}{0.95}
\renewcommand{\dbltopfraction}{0.95}
\renewcommand{\textfraction}{0.05}
\renewcommand{\floatpagefraction}{0.85}
\renewcommand{\dblfloatpagefraction}{0.85}
\setlength{\textfloatsep}{6pt plus 1pt minus 2pt}
\setlength{\floatsep}{6pt plus 1pt minus 2pt}
\setlength{\intextsep}{6pt plus 1pt minus 2pt}
\setlength{\dbltextfloatsep}{6pt plus 1pt minus 2pt}
\setlength{\dblfloatsep}{6pt plus 1pt minus 2pt}
\newcommand{\figref}[1]{Figure~\ref{#1}}
\newcommand{\figrefs}[2]{Figures~\ref{#1} and~\ref{#2}}
\newcommand{\figrangeref}[2]{Figures~\ref{#1}--\ref{#2}}
\newcommand{\tabref}[1]{Table~\ref{#1}}
\newcolumntype{L}[1]{>{\raggedright\arraybackslash}p{#1}}
\newcolumntype{R}[1]{>{\raggedleft\arraybackslash}p{#1}}
\newcolumntype{C}[1]{>{\centering\arraybackslash}p{#1}}
\newcounter{appendixsection}
\newcounter{appendixsubsection}[appendixsection]
\newcounter{appendixdomain}[appendixsubsection]
\renewcommand{\theappendixsection}{\Alph{appendixsection}}
\renewcommand{\theappendixsubsection}{\theappendixsection.\arabic{appendixsubsection}}
\renewcommand{\theappendixdomain}{\theappendixsubsection.\arabic{appendixdomain}}
\newcommand{\appendixsection}[2]{%
  \refstepcounter{appendixsection}%
  \section*{\theappendixsection. #1}%
  \label{#2}%
}
\newcommand{\appendixsubsection}[1]{%
  \refstepcounter{appendixsubsection}%
  \subsection*{\theappendixsubsection\ #1}%
}

\urlstyle{rm}
\def\UrlFont{\rm}

\title{Solver-Integrated Adversarial Attacking and Training of Neural Operators}

\author{Anonymous Authors}
\affiliations{Anonymous Affiliations}

\begin{document}

\maketitle
\begin{abstract}
Neural operators are widely used as fast surrogates for numerical PDE
solvers, mapping input functions to solution functions. However, their
generalizability and robustness are not yet clearly defined in the
operator-learning setting, which differs from traditional adversarial robustness
definitions. This paper studies the generalizability and robustness of a
learned neural operator from a solver-integrated perspective, addressing the
challenge that the output of a learned operator and a numerical solver tends to
change in tandem under input perturbation. First, we formalize the definition
of generalization and robustness through a model-solver error operator,
identifying fixed-input model-solver loss as generalization metric, and
norm-bounded adversarial attack loss increase and Jacobian-error function norm as
robustness metric. Second, we
identify the solver-integrated adversarial attack as appropriate
for PDE operator learning and show why model-only or fixed-ground-truth attacks
can be insufficient when the solver output also changes with the input. Third, we
develop solver-integrated adversarial training methods for neural operators.
Experiments on representative PDE benchmarks show that this solver-integrated
adversarial training clearly improves both generalizability and robustness.
Deeper solver integration yields more effective attacks, more informative
samples, and more efficient training than less integrated alternatives. More
broadly, this paper is a carefully chosen instance of a wider
adversarial-regression setting: the PDE solver provides a ground-truth oracle
that reveals how true outputs change under input perturbations and supplies
their input gradients, information unavailable in many prior problem settings.
\end{abstract}

\section{Introduction}

Neural operators, including Fourier Neural Operators and DeepONets, learn maps
between function spaces and are widely used as fast data-driven
surrogates for numerical PDE solvers \citep{li2021fourier,lu2021deeponet}.
They map input functions to solution functions in one forward pass, making them
attractive when many related PDE solution queries are needed. Practical use, however,
requires the learned operator to generalize beyond the training dataset and
remain close to the solver when the input function is perturbed. Recent neural
operator studies have therefore begun to examine adversarial robustness and
robustness-aware training
\citep{adesoji2022evaluating,huang2026stablepdenet,roy2026beyond}.

Prior work has often not clearly separated the notions of generalization and
robustness.
Classical theory defines generalization through expected-versus-empirical loss
\citep{bousquet2002stability,xu2012robustness}. Several modern robustness
notions are defined by worst-case loss increase under norm-bounded
perturbations \citep{madry2018towards}, and global Lipschitz constants and
local spectral norms are also used to measure robustness
\citep{cisse2017parseval,huang2026stablepdenet}. We inherit these ideas and
specialize them to the operator learning setting. Generalization is a static
definition: the norm of the output of our error operator. Robustness is defined
in a dynamic way: how the output of our error operator amplifies under an input
perturbation. We further develop Jacobian-error, which is cheaper to
compute and closely aligned with adversarial loss increase in the local
robustness setting.

Most classical adversarial attacks were developed for classification models, where it
is reasonable to assume that a small input perturbation should not change the
ground-truth label \citep{goodfellow2015explaining,madry2018towards}. However,
the PDE operator learning problem is not only a regression problem, but a
regression problem with an existing numerical solver as a trusted oracle. When
the input function is perturbed, the correct solution should change accordingly
but stay close to the solver's output. Therefore, the model is not expected to
be robust by itself; rather, the model-solver discrepancy is expected to be
robust. We therefore define the error operator as the discrepancy between the
learned operator and the solver.

Building on this view, we develop solver-integrated adversarial training. The
method is a teacher-student solver distillation procedure: the numerical solver
is the oracle teacher, and the neural operator is retrained on informative
inputs where the local model-solver discrepancy is large. This is the special
opportunity that makes the framework operational here: in ordinary regression or
vision problems, perturbing an input can also change the true output, yet the
true input-output map and its gradient are usually unavailable; in PDE operator
learning, the solver can relabel each perturbed input and can be differentiated
through. Experiments show that exploiting this oracle improves generalization
and robustness. More broadly, the paper should be read as a carefully chosen
instance of a broader adversarial-regression problem with access to a
ground-truth oracle, rather than as a claim limited to neural operators alone.

Our main contributions can be summarized as follows:
\begin{itemize}
  \item In the setting of operator learning, we propose a new robustness
  definition and metric.
  \item We identify the most accurate and appropriate target loss for
  neural-operator adversarial attack, which involves deep solver integration.
  \item We propose solver-integrated adversarial training for neural operators
  to improve both generalization and robustness.
\end{itemize}

\section{Related Work}

\noindent\textbf{Traditional adversarial attacks.}
Most classical adversarial attacks were designed for classification, where the
objective is usually to maximize a label loss such as
$\ell(f_\theta(x+\delta),y)$ under a perturbation constraint
\citep{goodfellow2015explaining,madry2018towards}. These attacks mainly differ
in the optimization procedure, perturbation constraint, sparsity or
structure, or available information, while retaining the assumption that the
ground-truth label remains fixed
\citep{carlini2017towards,papernot2016limitations,moosavi2016deepfool,brendel2017decision,chen2019hopskipjump,andriushchenko2019square,su2017onepixel,brown2017adversarialpatch,croce2020autoattack}.
Regression attacks and time-series forecasting attacks have also been studied
\citep{nguyen2018regression,tong2018adversarialregression,ribeiro2023linearregression,mode2020mtsregression,liu2022forecastingattacks}.
Their target losses use a fixed observed label or the model's own output. They
do not recompute the true loss with the true output for the perturbed input
inside the attack loop. They acknowledge that the true output should change
accordingly with input perturbations, but lack
the generating function in that setting. Therefore, they assume that small
perturbations do not meaningfully change the observed label. However, in our operator-learning
setting, the numerical solver is exactly this generating function, so the true
output and the corresponding true loss can be recomputed for each perturbed
input inside the loop.
Under the framework below, previous attack losses contain relatively little solver
integration and fall into the first two target-loss categories,
$\mathcal{L}_1$ and $\mathcal{L}_2$.

\noindent\textbf{Neural-operator adversarial attacks.}
A closer study evaluates adversarial robustness for operator learning and finds
that neural-operator error can grow significantly as the perturbation norm
increases \citep{adesoji2022evaluating}. Although solver outputs are used as
references, the solver is not truly integrated into the attack loop. In the
forward pass, the solver output for an updated perturbed input is approximated
by a nearest-neighbor ground-truth output from a precomputed dictionary; in the
backward pass, the solver gradient is ignored and only the neural-operator
gradient is used. Besides, this study focuses on adversarial robustness
evaluation for neural operators, but does not propose a method for improving the
model's adversarial robustness.

\noindent\textbf{PDE-problem adversarial training.}
WbAR combines adversarial training with PINNs by moving collocation points toward
larger PDE residuals \citep{shi2025wbar}; StablePDENet applies adversarial
training to neural operators with a physics loss
\citep{huang2026stablepdenet}; and RAMS systematizes this residual-gradient
sample movement for PINNs and operator learning \citep{ouyang2025rams}. These
methods do not use true solver outputs as supervised targets inside the
adversarial loop. Instead, they use physics loss as a surrogate loss. The
drawbacks are that this physics loss is less informative than direct solver
supervision and that, in many time-dependent neural operator settings, the model
outputs only one single frame of the final state, so the temporal derivatives
required by a physics loss are unavailable. Other
operator-learning works use heuristic, active, or generative sampling:
neural-operator digital-twin studies use gradient-free sparse attacks and
attack-driven active learning with input denoising
\citep{roy2026vulnerabilities,roy2026beyond}, while adversarial autoencoders,
GANO, and turbulent-flow generative models use distributional or generative
objectives \citep{enyeart2024adversarial,rahman2022gano,oommen2025turbulent}.
These sampling strategies can be useful, but they are heuristic or
distributional rather than chosen according to the current solver-model discrepancy, and
they do not use solver gradients to choose perturbation directions.

\section{Preliminaries}
\subsection{Problem Formulation}

Let $\mathcal{X}$ and $\mathcal{Y}$ be Hilbert function spaces equipped with
norms $\|\cdot\|_{\mathcal{X},p}$ and $\|\cdot\|_{\mathcal{Y},q}$, respectively,
where $p=q=2$ by default. An input is a function $x\in\mathcal{X}$, such as an
initial condition or coefficient field, and the output is a function
$y\in\mathcal{Y}$, such as a final condition or solution function. Throughout
the paper, the model is the surrogate neural
operator $\mathcal{F}_\theta$, while the solver is the oracle ground-truth
input-to-output mapping $\mathcal{G}$. We assume these operators are Frechet
differentiable at $x$. Their difference defines the
model-solver error operator $\mathcal{E}_\theta$:
\begin{equation*}
\begin{aligned}
  \mathcal{F}_\theta,\mathcal{G},\mathcal{E}_\theta
  &:\mathcal{X}\to\mathcal{Y},\
  \mathcal{E}_\theta:=\mathcal{F}_\theta-\mathcal{G},\\
  e_\theta[x]
  &:=
  \mathcal{E}_\theta(x)
  =
  \mathcal{F}_\theta(x)-\mathcal{G}(x),\\
  D\mathcal{F}_\theta[x],D\mathcal{G}[x]
  &:\mathcal{X}\to\mathcal{Y},\
  D\mathcal{E}_\theta[x]:=D\mathcal{F}_\theta[x]-D\mathcal{G}[x].
\end{aligned}
\end{equation*}
For a perturbation increment $\eta=t h\in\mathcal{X}$, let $h\in\mathcal{X}$ be an
input-space direction with $\|h\|_{\mathcal{X},p}=1$. We take $t\to0$, and hence
$\|\eta\|_{\mathcal{X},p}\to0$:
\begin{align*}
  \mathcal{F}_\theta(x+\eta)
  &=
  \mathcal{F}_\theta(x)+D\mathcal{F}_\theta[x](\eta)
  +o(\|\eta\|_{\mathcal{X},p}),
  \\
  \mathcal{G}(x+\eta)
  &=
  \mathcal{G}(x)+D\mathcal{G}[x](\eta)
  +o(\|\eta\|_{\mathcal{X},p}),
  \\
  e_\theta[x+\eta]
  &=
  e_\theta[x]+D\mathcal{E}_\theta[x](\eta)
  +o(\|\eta\|_{\mathcal{X},p}).
\end{align*}
Here $D\mathcal{F}_\theta[x](\eta)$, $D\mathcal{G}[x](\eta)$, and
$D\mathcal{E}_\theta[x](\eta)$ are the corresponding first-order output
variation functions.
We define the radius-$\varepsilon$ perturbation ball as the admissible
constraint set for the perturbed input in the function space $\mathcal{X}$:
\begin{equation*}
  \mathcal{B}_{p,\varepsilon}(x)
  =
  \{x+\delta:\delta\in\mathcal{X},
  \|\delta\|_{\mathcal{X},p}\le\varepsilon\}.
\end{equation*}

\subsection{Generalization and Robustness Metrics}

\paragraph{1. Generalization.}
Generalization is the static model-solver discrepancy at a fixed input
function:
\begin{equation*}
  \mathcal{L}_{\mathrm{err}}(x)
  =
  \frac{1}{2}\|e_\theta[x]\|_{\mathcal{Y},2}^{2}
  =
  \frac{1}{2}
  \|\mathcal{F}_\theta(x)-\mathcal{G}(x)\|_{\mathcal{Y},2}^{2}.
\end{equation*}

\paragraph{2. Robustness.}
Robustness asks how the model's error changes instead of how the model changes
when the input function is perturbed. We organize it into three local metrics.

We use $h$ for an infinitesimal direction. Along $h$,
\begin{equation*}
  e_\theta[x+t h]
  =
  e_\theta[x]+tD\mathcal{E}_\theta[x](h)+o(t),
  \qquad t\to 0.
\end{equation*}
The first variation of $\mathcal{L}_{\mathrm{err}}$ is
\par\smallskip
\noindent
\begin{tabular*}{\columnwidth}{@{}l@{\extracolsep{\fill}}r@{}}
$\displaystyle D\mathcal{L}_{\mathrm{err}}[x](h):=\left.\frac{d}{dt}\frac{1}{2}\langle e_\theta[x+t h],e_\theta[x+t h]\rangle_{\mathcal{Y}}\right|_{t=0}$ & \\
$\displaystyle =\langle e_\theta[x],D\mathcal{E}_\theta[x](h)\rangle_{\mathcal{Y}}\!=\!\langle D\mathcal{E}_\theta[x]^*e_\theta[x],h\rangle_{\mathcal{X}}\!=\!\langle g_\theta(x),h\rangle_{\mathcal{X}}$ &
\end{tabular*}
\par\smallskip
where $g_\theta(x):=D\mathcal{E}_\theta[x]^*e_\theta[x]\in\mathcal{X}$, and
$D\mathcal{E}_\theta[x]^*:\mathcal{Y}\to\mathcal{X}$ is the adjoint derivative.
After discretization this becomes the vector
$J_{\mathcal{E}}(x)^\top e_\theta(x)$.
This derivative is the first-order growth rate of the squared model-solver error
loss when $x$ is perturbed along $h$. It is the function-space inner product
between the Jacobian-error function $g_\theta(x)$ and the perturbation direction
$h$.

\smallskip
\noindent\textbf{Jacobian norm.}
The first robustness metric is the induced norm of the derivative of the
error operator. In a Lipschitz view, a regional robustness bound uses the
largest derivative norm over an input region; at a fixed input $x$, this reduces
to the pointwise induced norm of the Frechet derivative. The Frechet derivative
is defined from perturbations $x+t h$ as $t\to0$; after this infinitesimal scale
is factored out, the Jacobian norm maximizes over normalized input-space
directions $h$:
\begin{equation*}
  \|D\mathcal{E}_\theta[x]\|_{p\to q}
  =
  \sup_{\|h\|_{\mathcal{X},p}=1}
  \|D\mathcal{E}_\theta[x](h)\|_{\mathcal{Y},q}.
\end{equation*}
For a small perturbation $t h$, the linearized change is
$tD\mathcal{E}_\theta[x](h)$. Thus the Jacobian norm gives the first-order
growth rate of the nonlinear error operator at $x$; the maximizing direction
$h$ is the Jacobian norm growth direction. After
discretization, with
$J=J_{\mathcal{E}}(x)$, one common case is
$2\to2:\sigma_{\max}(J)$ with the leading right singular vector as direction.

This Jacobian norm differs from a model-only local stability diagnostic
such as the spectral norm \(\|D\mathcal{F}_\theta[x]\|\), which StablePDENet
reports for neural operators \citep{huang2026stablepdenet}. A large
\(\|D\mathcal{F}_\theta[x]\|\) does not by itself imply poor solver-consistent
robustness, because the true solver \(\mathcal{G}\) can also have
non-negligible input-output sensitivity; both outputs are expected to change
when the input changes. The desired local property is that the model and solver
change similarly, \(D\mathcal{F}_\theta[x]\approx D\mathcal{G}[x]\), so the
error derivative \(D\mathcal{E}_\theta[x]\) is small.

\smallskip
\noindent\textbf{Adversarial attack loss increase.}
For a fixed input function $x$, the attack optimizes over the perturbation
$\delta$; the model, solver, and clean input are fixed. This metric directly
measures the worst loss (error's norm) increase over a norm-bounded perturbation
ball, rather than a local error change's norm. For one perturbation, define
\begin{align*}
  \Delta\mathcal{L}_{\mathrm{err}}(\delta)
  &:=
  \mathcal{L}_{\mathrm{err}}(x+\delta)
  -
  \mathcal{L}_{\mathrm{err}}(x).
\end{align*}
Since the clean error $\|e_\theta[x]\|_{\mathcal{Y}}^2$ is constant in
$\delta$, maximizing $\|e_\theta[x+\delta]\|_{\mathcal{Y}}^2$ and maximizing
$\Delta\mathcal{L}_{\mathrm{err}}(\delta)$ give the same perturbation. We
choose a budget-feasible $\delta$ that maximizes this loss increase; the chosen
$\delta$ is the finite-budget adversarial perturbation direction.

\smallskip
\noindent\textbf{Jacobian-error.}
The third robustness quantity is Jacobian-error, defined by the
input-space function
$g_\theta(x)=D\mathcal{E}_\theta[x]^*e_\theta[x]$. It combines the derivative
of the error operator with the current error field. Its scalar growth value is
$\|g_\theta(x)\|_{\mathcal{X}}$, and its growth direction is
$g_\theta(x)/\|g_\theta(x)\|_{\mathcal{X}}$ when $g_\theta(x)\neq0$.

\smallskip
\noindent\textbf{Comparison of Jacobian-error and finite-budget attack.}
In the small-budget limit, the finite-budget attack reduces to
Jacobian-error. Write
$\delta=\varepsilon h$ with $\|h\|_{\mathcal{X}}\le 1$. As
$\varepsilon\to0$, the first-order expansion is uniform over
$\|h\|_{\mathcal{X}}\le1$, so
\begin{align*}
  \Delta\mathcal{L}_{\mathrm{err}}^*(\varepsilon)
  &=
  \max_{\|h\|_{\mathcal{X}}\le1}
  \Delta\mathcal{L}_{\mathrm{err}}(\varepsilon h)
  \approx
  \varepsilon\max_{\|h\|_{\mathcal{X}}\le 1}
  \langle g_\theta(x),h\rangle_{\mathcal{X}}
  \\
  &=
  \varepsilon\|g_\theta(x)\|_{\mathcal{X}}
  \quad
  \left(
  h^*=g_\theta(x)/\|g_\theta(x)\|_{\mathcal{X}}
  \right).
\end{align*}
Thus, in the small-budget limit, the finite-budget adversarial attack has the
same first-order growth value and growth direction as Jacobian-error.

\smallskip
\noindent\textbf{Comparison of Jacobian-error and Jacobian norm.}
Jacobian-error differs from the Jacobian norm because the two quantities
optimize different local objectives. Jacobian-error measures the
error-aligned first-order growth of the
final squared error:
\begin{align*}
  \frac{1}{2}\|e_\theta[x+h]\|_{\mathcal{Y}}^2
  -
  \frac{1}{2}\|e_\theta[x]\|_{\mathcal{Y}}^2
  &\approx
  \bigl\langle e_\theta[x],
  D\mathcal{E}_\theta[x](h)\bigr\rangle_{\mathcal{Y}}
\end{align*}
This first-order term depends on both the current error $e_\theta[x]$ and the
local Jacobian through $D\mathcal{E}_\theta[x](h)$. The Jacobian norm measures
the largest local displacement of the error field:
\begin{equation*}
  \|e_\theta[x+h]-e_\theta[x]\|_{\mathcal{Y}}
  \approx
  \|D\mathcal{E}_\theta[x](h)\|_{\mathcal{Y}}
\end{equation*}
Thus, the Jacobian norm maximizes the local error displacement norm, not the
error norm (loss) increase. Consequently, in this setting the Jacobian norm is
not the most suitable robustness proxy; the Jacobian-error norm is more
directly aligned with finite adversarial attack loss increase. For the $2\to2$
norm case ($p=q=2$),
the Jacobian-norm direction is the leading right singular direction of
$D\mathcal{E}_\theta[x]$, governed by
$D\mathcal{E}_\theta[x]^*D\mathcal{E}_\theta[x]$. Empirically, Jacobian-error
tracks finite adversarial attack loss increase more closely than Jacobian norm
for small $\varepsilon$. It is also less
computationally intensive: for a discretized
$J_{\mathcal{E}}\in\mathbb{R}^{m\times n}$ and error vector
$e_\theta\in\mathbb{R}^{m}$, computing
$J_{\mathcal{E}}^\top e_\theta$ and its norm costs $O(mn)$, whereas computing
the Jacobian norm through the leading singular value/vector costs
$O(mn\min\{m,n\})$ for an exact dense SVD or $O(kmn)$ for $k$ iterative
matrix-vector passes.

\section{Proposed Framework}

\subsection{Solver-Integrated Adversarial Attack}

The attack optimizes the chosen loss over the finite perturbation $\delta$,
with $x$ fixed. For iterative variants, $x_k=x+\delta_k$ is the perturbed input
at step $k$, and $\operatorname{sg}(\cdot)$ denotes stop-gradient. We use
three primary losses and several implementation variants:
\begin{equation*}
\begin{aligned}
\mathcal{L}_1(\delta)
&=\|\mathcal{F}_\theta(x+\delta)-\mathcal{F}_\theta(x)\|_{\mathcal{Y},q},\\
\mathcal{L}_2(\delta)
&=\|\mathcal{F}_\theta(x+\delta)-\mathcal{G}(x)\|_{\mathcal{Y},q},\\
\mathcal{L}_3(\delta)
&=\|\mathcal{F}_\theta(x+\delta)-\mathcal{G}(x+\delta)\|_{\mathcal{Y},q}=\|e_\theta[x+\delta]\|_{\mathcal{Y},q},\\
\mathcal{L}_2^{\mathrm{fixed}}(\delta_k)
&=\|\mathcal{F}_\theta(x_k)-\mathcal{G}(x)\|_{\mathcal{Y},q},\\
\mathcal{L}_3^{\mathrm{full}}(\delta_k)
&=\|\mathcal{F}_\theta(x_k)-\mathcal{G}(x_k)\|_{\mathcal{Y},q},\\
\mathcal{L}_{2,N}^{\mathrm{dict}}(\delta_k)
&=\|\mathcal{F}_\theta(x_k)-y_{\mathrm{dict},N}(x_k)\|_{\mathcal{Y},q},\\
\mathcal{L}_3^{\mathrm{sg}}(\delta_k)
&=\|\mathcal{F}_\theta(x_k)-\operatorname{sg}(\mathcal{G}(x_k))\|_{\mathcal{Y},q},\\
\mathcal{L}_{\mathrm{phys}}(\delta_k)
&=\lambda_{\mathrm{PDE}}\|R_{\mathrm{PDE}}(\mathcal{F}_\theta(x_k))\|^2\\
&\quad+\lambda_{\mathrm{BC}}\|R_{\mathrm{BC}}(\mathcal{F}_\theta(x_k))\|^2+\lambda_{\mathrm{IC}}\|R_{\mathrm{IC}}(x_k)\|^2.
\end{aligned}
\end{equation*}
For the squared-loss version, the stop-gradient and full $\mathcal{L}_3$
variants have different backpropagated first-order directions:
\begin{align*}
  \nabla_{\delta_k}\tfrac{1}{2}\mathcal{L}_3^{\mathrm{sg}}(\delta_k)^2
  &=
  D\mathcal{F}_\theta[x_k]^*
  (\mathcal{F}_\theta(x_k)-\mathcal{G}(x_k)),
  \nonumber\\
  \nabla_{\delta_k}\tfrac{1}{2}\mathcal{L}_3^{\mathrm{full}}(\delta_k)^2
  &=
  (D\mathcal{F}_\theta[x_k]-D\mathcal{G}[x_k])^*
  (\mathcal{F}_\theta(x_k)-\mathcal{G}(x_k)).
\end{align*}

Model-output-change objectives, including some regression attacks and
regularizers that compare perturbed and clean predictions, follow
$\mathcal{L}_1$
\citep{nguyen2018regression,mode2020mtsregression,liu2022forecastingattacks}.
Fixed-label classification, fixed-response regression, specified-target
regression/forecasting, and neural-operator tests with held or precomputed
solver targets follow $\mathcal{L}_2$
\citep{goodfellow2015explaining,madry2018towards,tong2018adversarialregression,ribeiro2023linearregression,liu2022forecastingattacks,adesoji2022evaluating}.
We are not aware of prior attack losses that optimize the full
$\mathcal{L}_3$.

$\mathcal{L}_{2,N}^{\mathrm{dict}}$ is an $\mathcal{L}_2$-like variant: it
replaces the online solver output by the precomputed solver output attached to
the nearest dictionary input, rather than calling the solver during the attack.
This is the structure used by dictionary-based FNO robustness evaluation
\citep{adesoji2022evaluating}. The stop-gradient variant uses the solver output
in the forward target but blocks the solver branch in backpropagation, whereas
$\mathcal{L}_3^{\mathrm{full}}$ differentiates through the model--solver
difference as shown above. Physics-residual attacks such as WbAR, StablePDENet,
and RAMS instead use $\mathcal{L}_{\mathrm{phys}}$: they do not compare with
solver outputs, but replace solver supervision by PDE, boundary, and
initial-condition residuals
\citep{shi2025wbar,huang2026stablepdenet,ouyang2025rams}.
Overall, these attack variants increase solver integration in the order
$\mathcal{L}_{\mathrm{phys}}$, $\mathcal{L}_1$, $\mathcal{L}_2$,
$\mathcal{L}_{2,N}^{\mathrm{dict}}$, $\mathcal{L}_3^{\mathrm{sg}}$, and
$\mathcal{L}_3^{\mathrm{full}}$.
Our experiments below show that deeper solver integration generally yields
stronger attacks and more effective adversarial training.

\subsection{Optimization Methods for Solver-Integrated Attacks}

We propose four optimizer variants, and many existing first-order attack
optimizers can be viewed within this framework. For a selected attack loss
$\mathcal{L}(\delta)$, let
$g_k=\nabla_\delta\mathcal{L}(\delta_k)$ denote the input-space gradient at
attack step $k$. Raw updates use $g_k$ directly. The steepest direction is
\begin{equation*}
  s_k \in \operatorname*{arg\,max}_{\|v\|_{\mathcal{X},p}\le 1}
  \langle g_k,v\rangle .
\end{equation*}

\[
\begin{array}{@{}lcc@{}}
\hline
& \textbf{Add} & \textbf{Replace} \\
\hline
\textbf{Raw}
&
\delta_{k+1}
=
\Pi_{\mathcal{B}_\varepsilon}
\left(\delta_k+\alpha g_k\right)
&
\delta_{k+1}
=
\varepsilon
\frac{g_k}{\|g_k\|_{\mathcal{X},p}}
\\[0.5em]
\textbf{Steepest}
&
\delta_{k+1}
=
\Pi_{\mathcal{B}_\varepsilon}
\left(\delta_k+\alpha s_k\right)
&
\delta_{k+1}
=
\varepsilon s_k
\\
\hline
\end{array}
\]

Raw add is standard PGD, and steepest add is $\ell_p$-steepest PGD.
Add-style methods accumulate directions and project back
to the perturbation ball, so early steps may move toward the boundary.
Replace-style methods discard the previous perturbation and directly place the
new perturbation on the boundary. Raw replace is similar to generalized power
iteration for solving Rayleigh quotient problems.
In our experiments, replace-style updates are fast and often strong on Burgers
and Darcy Flow, whereas add-style updates are more stable for recurrent
Navier--Stokes.

\subsection{Solver-Integrated Adversarial Training}

The training stage wraps the attack objectives above in an online
attack-then-train loop. We start from a trained baseline neural operator. At
step $t$, the current model is frozen during attack generation, and an attacked
batch $x^{\mathrm{adv}}=x+\delta_N$ is obtained by maximizing one of
$\mathcal{L}_1$, $\mathcal{L}_2$, $\mathcal{L}_3$, or
$\mathcal{L}_{\mathrm{phys}}$ for $N$ projected steps. This attack step searches
for locally high-loss input functions where the current model has large
model-solver discrepancy or large physics residual. The trainable model is
then updated on the attacked batch, and the updated weights become the next
current model.

This process can be summarized as a function-space min-max risk.  Let
$\mu_{\mathcal X}$ be the data distribution over the input function space
$\mathcal X$, so that $x\sim\mu_{\mathcal X}$.  The usual expected risk is
\(\mathcal R(\theta)=
\mathbb E_{x\sim\mu_{\mathcal X}}[\mathcal L_\theta(x)]\). The adversarial
attack moves the sample \(x\) toward a higher-loss
region using a perturbation \(\delta\) within the \(\varepsilon\)-ball. The
adversarial training objectives are therefore
\[
  \min_\theta \mathcal R_{\mathrm{robust}}(\theta)
  =
  \min_\theta
  \mathbb E_{x\sim\mu_{\mathcal X}}
  \left[
    \max_{\|\delta\|_{\mathcal X,p}\le \varepsilon}
    \mathcal L_\theta(x;\delta)
  \right].
\]

In our experiments, we also include two hand-designed random perturbation
training methods. These variants draw perturbations from a fixed random sampler
with $\varepsilon$ jitter that is independent of the model at the current
training step. They differ only in target: \emph{random solver} uses
$y=\mathcal{G}(x+\delta)$, while \emph{random clean} keeps
$y=\mathcal{G}(x)$.

\section{Experiments}

\subsection{Benchmark PDEs}

We evaluate on the three canonical PDE operator-learning benchmarks used in the
Fourier Neural Operator (FNO) paper \citep{li2021fourier}: 1D Burgers, 2D
Darcy Flow, and 2D Navier--Stokes. Due to page limits, many experimental
results have been moved to the supplement.

\subsection{Solver-Integrated Adversarial Attacks}
\label{subsec:exp-solver-integrated-attacks}

Except for Burgers at \(\nu=10^{-2}\), where we also include DeepONet, all
other experiments use FNO. For Burgers and recurrent Navier--Stokes, we compare
\(\mathcal L_1\), \(\mathcal L_2\), and \(\mathcal L_3\); for Darcy Flow, we
additionally compare \(\mathcal L_{\mathrm{phys}}\). In visualizations and
auxiliary sweeps, we also compare \(\mathcal L^{\mathrm{sg}}_3\),
\(\mathcal L^{\mathrm{dict}}_{2,N}\) with \(N\in\{200,2000,20000\}\), and
\(\mathcal L^{\mathrm{fixed}}_2\). Across these settings, the
attack-loss increase grows approximately as a power law in the attack budget,
$\Delta\mathcal{L}\approx C\varepsilon^a$, as summarized in
\figref{fig:main-cross-benchmark-epsilon-loss}.

\setcounter{figure}{0}
\begin{figure}[t]
  \centering
  \includegraphics[width=\columnwidth]{../figures/adversarial_attack_appendix/attack/cross_benchmark/epsilon_sweeps/epsilon_vs_loss_delta_best_all_loglog_powerlaw_annotated_1p3x_mathtext.png}
  \caption{Power-law relationship between attack budget and loss increase}
  \label{fig:main-cross-benchmark-epsilon-loss}
\end{figure}

\begin{figure}[t]
  \centering
  \includegraphics[width=\columnwidth,keepaspectratio]{../selected_figure_exports/fig10_crop_first2rows_first2cols.png}\\
  \includegraphics[width=\columnwidth,keepaspectratio]{../selected_figure_exports/fig10_crop_first2rows_last2cols.png}\\
  \includegraphics[width=\columnwidth,keepaspectratio]{../selected_figure_exports/fig10_crop_thirdrow_first2cols.png}
  \caption{Burgers FNO \(L_2\) PGD attack variants.}
  \label{fig:main-burgers-attack-variant-groups}
\end{figure}

\begin{table}[t]
\centering
\normalsize
\setlength{\tabcolsep}{1.1pt}
\textbf{Burgers} (\(100\) attack steps, \(N=50\) samples)\\[0.15em]
\begin{tabular}{@{}lccccc@{}}
\toprule
Budget & Step size & \(\mathcal L_1 \uparrow\) & \(\mathcal L_2 \uparrow\) &
\(\boldsymbol{\mathcal L}_3 \uparrow\) & \(\mathcal L_3\) win \(\uparrow\) \\
\midrule
\(\varepsilon=0.125\) & 0.0125 & 0.1027 & 0.1021 & \textbf{0.1306}\(^{*}\) & 50/50 \\
\(\varepsilon=0.25\) & 0.025 & 0.1008 & 0.0999 & \textbf{0.1664}\(^{*}\) & 50/50 \\
\(\varepsilon=0.5\) & 0.05 & 0.0982 & 0.0963 & \textbf{0.2793}\(^{*}\) & 50/50 \\
\(\varepsilon=0.75\) & 0.075 & 0.0976 & 0.0934 & \textbf{0.4612}\(^{*}\) & 50/50 \\
\(\varepsilon=1.0\) & 0.1 & 0.0988 & 0.0913 & \textbf{0.7262}\(^{*}\) & 50/50 \\
\bottomrule
\end{tabular}

\vspace{0.1em}
\setlength{\tabcolsep}{0.6pt}
\textbf{Darcy Flow} (\(100\) attack steps, \(N=20\) samples)\\[0.15em]
\begin{tabular}{@{}lcccccc@{}}
\toprule
Budget & Step size & \(\mathcal L_1 \uparrow\) & \(\mathcal L_2 \uparrow\) &
\(\boldsymbol{\mathcal L}_3 \uparrow\) & \(\mathcal L_{\mathrm{phys}} \uparrow\) &
\(\mathcal L_3\) win \(\uparrow\) \\
\midrule
$K=100$ & 5 & 0.0254 & 0.0256 & \textbf{0.0307}\(^{*}\) & 0.0245 & 14/20 \\
$K=250$ & 5 & 0.0266 & 0.0270 & \textbf{0.0403}\(^{*}\) & 0.0245 & 17/20 \\
$K=437$ & 5 & 0.0276 & 0.0284 & \textbf{0.0500}\(^{*}\) & 0.0244 & 17/20 \\
$K=875$ & 5 & 0.0288 & 0.0309 & \textbf{0.0771}\(^{*}\) & 0.0245 & 20/20 \\
\bottomrule
\end{tabular}

\vspace{0.1em}
\setlength{\tabcolsep}{1.1pt}
\textbf{Navier--Stokes} (\(50\) attack steps, \(N=20\) samples)\\[0.15em]
\begin{tabular}{@{}lccccc@{}}
\toprule
Budget & Step size & \(\mathcal L_1 \uparrow\) & \(\mathcal L_2 \uparrow\) &
\(\boldsymbol{\mathcal L}_3 \uparrow\) & \(\mathcal L_3\) win \(\uparrow\) \\
\midrule
\(\varepsilon=8\) & 0.25 & 64.979 & 59.567 & \textbf{126.832}\(^{*}\) & 20/20 \\
\(\varepsilon=16\) & 0.5 & 68.913 & 61.469 & \textbf{186.683}\(^{*}\) & 20/20 \\
\(\varepsilon=32\) & 1 & 118.588 & 69.505 & \textbf{276.760}\(^{*}\) & 20/20 \\
\(\varepsilon=64\) & 2 & 190.810 & 88.554 & \textbf{324.618}\(^{*}\) & 18/20 \\
\bottomrule
\end{tabular}
\caption{Attack Objective Comparison. \(\mathcal L_3\) is used as the
true loss here, reported as the sample mean. Superscript \(^{*}\) marks a
one-sided test against the second-highest objective (\(q\le 0.05\)).}
\label{tab:attack_objective_true_l3_summary}
\end{table}

The main result is simple: to make the model output differ strongly from the
solver output, the attack should optimize $\mathcal{L}^{\mathrm{full}}_3$
directly. As shown in \tabref{tab:attack_objective_true_l3_summary}, attacks based on
$\mathcal{L}_1$, $\mathcal{L}_2$, $\mathcal{L}_{\mathrm{phys}}$,
$\mathcal{L}^{\mathrm{fixed}}_2$, $\mathcal{L}^{\mathrm{dict}}_{2,N}$, or
$\mathcal{L}^{\mathrm{sg}}_3$ increase the true $\mathcal{L}_3$ more
slowly and reach a smaller final $\mathcal{L}_3$. Therefore, we need to use the
solver both in the forward pass and in the backward pass through
$D\mathcal{G}$. See Figures~\ref{fig:main-burgers-attack-variant-groups}
and~\ref{fig:main-ns-steepest-add-final} for attack visualizations.

Our second finding concerns frequency content. In Burgers and Navier--Stokes,
perturbations generated by full $\mathcal{L}_3$ concentrate on higher-frequency
modes than those from other objectives. The local SVD diagnostic quantifies the
mechanism: for Burgers FNO at \(\nu=10^{-2}\), after removing the DC mode, the
\(q_{50}\) frequency of the rank-1 right singular vector is \(0.20\%\) of the
Fourier grid for both \(D\mathcal{F}_\theta[x]\) and \(D\mathcal{G}[x]\), but
\(15.7\%\) for \(D\mathcal{E}_\theta[x]\). For ranks 2--4, the corresponding
\(q_{50}\) values remain at most \(0.39\%\) for \(D\mathcal{F}_\theta[x]\) and
\(D\mathcal{G}[x]\), while \(D\mathcal{E}_\theta[x]\) ranges from \(45.4\%\) to
\(69.3\%\).
Thus, to maximize the model--solver difference, the perturbation must move
along this sharper high-frequency direction.

\setcounter{figure}{2}
\begin{figure*}[t]
  \centering
  \includegraphics[width=\textwidth,keepaspectratio]{../selected_figure_exports/fig26_composite_ns_loss3_with_loss_curve_spectrum.png}
  \caption{Loss 3 creates the largest model-solver discrepancy for 2D
  Navier--Stokes. Bottom: all-sample mean true-loss curve and perturbation
  spectrum.}
  \label{fig:main-ns-steepest-add-final}
\end{figure*}

\subsection{Projected Attack Optimization Methods}
\label{subsec:exp-projected-optimization}

The optimizer ablations give a surprising pattern: for Burgers and Darcy Flow,
replace-style updates move to the perturbation boundary on the first
replacement update and then attain high final true \(\mathcal L_3\). In the
saved Burgers diagnostic, steepest replace ends at \(\mathcal L_3=7.069\),
while steepest add ends at \(6.339\), a \(10.3\%\) gap relative to replace. In
Darcy Flow, binary replace ends at \(0.05285\), while binary add ends at
\(0.04809\), a \(9.0\%\) gap relative to replace. Mechanism diagnostics suggest
two reasons: replace reaches the constraint boundary immediately instead of
slowly moving outward, and the Burgers/Darcy loss landscapes contain broad
connected high-loss regions, so a coarse boundary move can already enter such a
region.
Recurrent Navier--Stokes is different: add-style optimization reaches a much
larger \(\mathcal L_3\) than replace, indicating a narrower and sharper
high-loss region that is hard to hit without more refined iterative
optimization.

\subsection{Adversarial Training}
\label{subsec:exp-adversarial-training}

Because different training methods have different per-epoch time costs, we
compare all models at a similar total wall-clock training time. We measure both
generalization to shifted datasets and robustness under adversarial attack.
These broad generalization datasets are generated by shifting the input
distribution: we vary GRF kernel families, correlation lengths, and
spectral-decay parameters, include non-GRF function families, and change
amplitude ranges or thresholding transforms.
For Burgers and Darcy Flow, we conduct detailed and complete experimental
comparisons across various training methods.

\paragraph{Generalization.}
On Burgers (Figure~\ref{fig:main-burgers-generalization-top-row}), Loss 3 gives
the lowest generalization relative \(L_2\) on \(48\) of the \(50\) shifted datasets.
Random solver is the second-best method and performs very well early in
wall-clock time, but after about four hours it shows clear overfitting: its
generalization loss starts to increase. We infer that the main reason is that
its \(\delta\) is manually prescribed in advance from a fixed statistical
distribution, so it cannot use adversarial attack against the current model to
generate perturbations along the model's current weak directions. As shown in
Figure~\ref{fig:main-burgers-delta-frequency-line-plot}, \(\mathcal L_3\)
adversarial training generates perturbations that clearly become
higher-frequency as training time increases, while \(\mathcal L_1\) and
\(\mathcal L_2\) training do not show this obvious high-frequency shift. Darcy
Flow (Figure~\ref{fig:main-darcy-generalization-top-row}) shows the same
pattern, with Loss 3 giving the lowest
generalization relative \(L_2\) on \(47\) of the \(50\) shifted datasets. For
both Burgers and Darcy Flow, the
\(\mathcal L_3\)-trained model is not the best model on the original train/test
datasets. This suggests a possible trade-off between in-distribution and
out-of-distribution performance: a model that performs better on shifted
datasets can perform worse on the original train/test distribution.

\setcounter{figure}{3}
\begin{figure*}[t]
  \centering
  \includegraphics[width=\textwidth,keepaspectratio]{../selected_figure_exports/fig38_crop_top1row.png}
  \caption{Burgers selected shifted-dataset generalization loss curves of all
  training methods. (Random clean is omitted here because it performs poorly
  and appears outside the plotted loss scale.)}
  \label{fig:main-burgers-generalization-top-row}
  \includegraphics[width=\textwidth,keepaspectratio]{../selected_figure_exports/fig47_crop_top1row.png}
  \caption{Darcy Flow selected shifted-dataset generalization loss curves of all training methods.}
  \label{fig:main-darcy-generalization-top-row}
\end{figure*}

\begin{figure}[t]
  \centering
  \includegraphics[width=\columnwidth,keepaspectratio]{../selected_figure_exports/fig41_crop_line_plot.png}
  \caption{Increasingly higher-frequency perturbation in Burgers Loss-3
  adversarial training.}
  \label{fig:main-burgers-delta-frequency-line-plot}
\end{figure}

For 2D Navier--Stokes, we do not compare different training methods; we only
run \(\mathcal L_3\) adversarial training, with six rounds in total. The main
reason is that Navier--Stokes adversarial training requires very long
wall-clock time and very large GPU memory, and we do not have sufficient
computational resources. After training, 364 out of 402 datasets show lower
MSE.

\paragraph{Robustness.}
\begin{table}[t]
\centering
\normalsize
\setlength{\tabcolsep}{5pt}
\textbf{Burgers}\\[0.2em]
\setlength{\tabcolsep}{1.4pt}
\begin{tabular}{@{}lcccccc@{}}
\toprule
Metric & Base & \(\mathcal L_1\) & \(\mathcal L_2\) &
\(\mathcal L_3\) & \shortstack{Rand.\\clean} & \shortstack{Rand.\\solver} \\
\midrule
Attack \(\Delta L\downarrow\) &
0.0092 & 0.0058 & 0.0056 & \textbf{0.0031}\(^{*}\) & 0.0428 & 0.0080 \\
\(J\)-error \(\downarrow\) &
0.4630 & 0.3159 & 0.3433 & \textbf{0.1119}\(^{*}\) & 4.7850 & 0.3712 \\
\bottomrule
\end{tabular}

\vspace{0.05em}
\textbf{Darcy Flow}\\[0.2em]
\setlength{\tabcolsep}{3pt}
\begin{tabular}{@{}lcccc@{}}
\toprule
Metric & Base & \(\mathcal L_1\) & \(\mathcal L_2\) &
\(\mathcal L_3\) \\
\midrule
Attack \(\Delta L\downarrow\) &
9.681e-6 & 3.927e-6 & 3.565e-6 & \textbf{1.699e-6}\(^{*}\) \\
\(J\)-error \(\downarrow\) &
9.293e-5 & 9.922e-5 & 9.225e-5 & \textbf{6.157e-5}\(^{*}\) \\
\bottomrule
\end{tabular}

\vspace{0.05em}
\begin{tabular}{@{}lccc@{}}
\toprule
Metric & \(\mathcal L_{\mathrm{phys}}\) &
\shortstack{Rand.\\clean} & \shortstack{Rand.\\solver} \\
\midrule
Attack \(\Delta L\downarrow\) & 3.974e-6 & 3.691e-6 & 4.824e-6 \\
\(J\)-error \(\downarrow\) & 1.034e-4 & 9.112e-5 & 1.236e-4 \\
\bottomrule
\end{tabular}
\caption{Adversarial trained models robustness. Values are sample means.
Attack loss uses \(50\) shifted datasets with \(50\)
samples in each dataset.  The Jacobian-error metric uses \(25\) selected
data points.  Superscript \(^{*}\) marks a one-sided test against the
second-lowest value (\(q\le 0.05\)).}
\label{tab:adversarial_training_robustness_summary}
\end{table}

\tabref{tab:adversarial_training_robustness_summary} reports finite-budget
attack-loss increase and the Jacobian-error norm
\(\|J_{\mathcal E}^{\top}e_\theta\|_2\). Loss 3 is strongest on both Burgers
and Darcy Flow for attack loss and Jacobian-error robustness. We do not report
the Jacobian spectral norm because it correlates less with finite-budget
attack-loss increase and perturbation directions, especially at small
\(\varepsilon\).

\section{Conclusion}

This paper studies the generalizability and robustness of neural operators from
a solver-integrated perspective. We define both quantities through
model--solver discrepancy under distribution shift and adversarial
perturbation, show why model-only robustness indicators are insufficient for PDE
operator learning. We then propose a solver-integrated adversarial attack whose
forward loss and backward gradient both use the PDE solver, so the attack
directly maximizes the relevant model--solver error. From this attack, we derive
an adversarial-training procedure that generates solver-relabelled adversarial
batches online. Experiments on Burgers, Darcy Flow, and Navier--Stokes
settings show that this solver-integrated adversarial training clearly improves
both generalizability and robustness over non-solver-integrated objectives,
without manually selecting special shifted training datasets.
More broadly, the framework is not limited to operator learning: any
neural-network system with access to a ground-truth oracle can adopt this
formulation for adversarial attacks and adversarial training.

The main limitation is computational cost. Solver-integrated attacks and
training require expensive solver forward passes and gradients through the
solver. For time-dependent PDEs, the solver backward pass must keep the
intermediate states and stage variables from many time steps in GPU memory so
that gradients can be propagated back through the entire rollout.
This creates a large memory footprint and a long wall-clock time cost, especially for
NS.
\bibliography{references}

\end{document}

